% ======================================================================
% Appendix A. Identity term and effective volume
% PATH: src/appendices/A-identity-term-effective-volume.tex
% ======================================================================

\appendix
\chapter{The identity term and effective volume computations}
\label{app:A-identity-effective-volume}

\section{Purpose and scope}
\label{sec:A.scope}

The purpose of this appendix is to record, in a self-contained form, the computation of the identity contribution
appearing in the localized trace formula, together with the effective volume constants arising from the local Weyl
law regime.  These computations are standard in principle, but the localized framework requires a precise accounting
of normalizations and the dependence on the cutoff profile.

Throughout, we work in the setting of a finite-area hyperbolic surface
\[
M = \Gamma \backslash \mathbb{H},
\]
equipped with the hyperbolic metric of curvature $-1$, and we use the spectral conventions fixed globally in
Chapter~\ref{chap:02-preliminaries} and Appendix~\ref{app:G-normalizations}.

\section{The wave kernel and the identity contribution}
\label{sec:A.wave-identity}

Let $\Delta$ denote the (positive) Laplace--Beltrami operator on $M$.  For a test function $h$ belonging to the
admissible class defined in Chapter~\ref{chap:07-localized-trace}, we consider the localized spectral operator
\[
h(\sqrt{\Delta}-\lambda),
\qquad \lambda\ge 1,
\]
defined by functional calculus.  Its Schwartz kernel may be written in terms of the wave group as
\begin{equation}
\label{eq:A.functional-calculus}
h(\sqrt{\Delta}-\lambda)
=
\frac{1}{2\pi}\int_{\mathbb{R}} \widehat{h}(t)\,e^{-it\lambda}\,e^{it\sqrt{\Delta}}\,dt,
\end{equation}
where $\widehat{h}$ is the Fourier transform of $h$ under the normalization fixed in
Appendix~\ref{app:G-normalizations}.

The trace of $h(\sqrt{\Delta}-\lambda)$ is given by integrating the kernel on the diagonal.  In the geometric
expansion underlying the localized trace formula, the identity contribution corresponds to the $\gamma = e$
summand in the group trace decomposition.  Equivalently, it is the contribution of the diagonal singularity of the
wave kernel on the universal cover.

\subsection{Local parametrix near the diagonal}
\label{subsec:A.parametrix}

Let $U(t;z,w)$ denote the wave kernel on $\mathbb{H}$:
\[
U(t;z,w) = \bigl(e^{it\sqrt{\Delta_{\mathbb{H}}}}\bigr)(z,w).
\]
The kernel is a function of the hyperbolic distance $r=d(z,w)$ and admits a Hadamard-type parametrix near $r=0$.
In particular, for $|t|$ sufficiently small, the diagonal behavior of $U(t;z,z)$ has an expansion of the form
\begin{equation}
\label{eq:A.wave-diagonal-expansion}
U(t;z,z)
=
\frac{1}{2\pi}\,\frac{\cosh(t/2)}{\sinh(t/2)}
+
R(t;z),
\end{equation}
where $R(t;z)$ is smooth in $t$ near $t=0$ and bounded uniformly in $z$ on compacta.  The singular term in
\eqref{eq:A.wave-diagonal-expansion} is universal and encodes the local Weyl density.

For the purposes of the trace formula, it is convenient to isolate the universal singular part by defining the
identity distribution
\begin{equation}
\label{eq:A.identity-distribution}
\mathcal{I}[h]
:=
\frac{\vol(M)}{2\pi}\int_{\mathbb{R}} h(r)\,r\,\tanh(\pi r)\,dr,
\end{equation}
which is the classical identity term in the Selberg trace formula.

The main task of this appendix is to explain how \eqref{eq:A.identity-distribution} reduces, under localization, to
the effective volume constant appearing in the local Weyl law.

\section{The Plancherel density and the local Weyl constant}
\label{sec:A.plancherel}

Let $\mu_{\mathrm{Pl}}(r)$ denote the Plancherel density for $\mathbb{H}$:
\begin{equation}
\label{eq:A.plancherel-density}
\mu_{\mathrm{Pl}}(r)
=
\frac{1}{2\pi}\,r\,\tanh(\pi r),
\qquad r\in\mathbb{R}.
\end{equation}
Then the identity term \eqref{eq:A.identity-distribution} may be rewritten as
\begin{equation}
\label{eq:A.identity-plancherel}
\mathcal{I}[h]
=
\vol(M)\int_{\mathbb{R}} h(r)\,\mu_{\mathrm{Pl}}(r)\,dr.
\end{equation}

We recall the asymptotic behavior
\begin{equation}
\label{eq:A.plancherel-asymptotic}
\mu_{\mathrm{Pl}}(r)
=
\frac{1}{2\pi}\,r
+
O\!\bigl(r e^{-2\pi |r|}\bigr),
\qquad |r|\to\infty,
\end{equation}
which shows that, at large spectral parameter, the hyperbolic Plancherel density is asymptotic to the Euclidean
density $\frac{1}{2\pi}r$.

\subsection{Localized cutoff profiles}
\label{subsec:A.localized-cutoff}

Let $\chi\in \mathcal{S}(\mathbb{R})$ be even with $\chi(0)=1$.  For $\eta>0$ define
\begin{equation}
\label{eq:A.chi-eta}
\chi_\eta(u)
:=
\chi(u/\eta).
\end{equation}
Fix $\lambda\ge 1$.  The localized spectral window is implemented by the weight $\chi_\eta(r-\lambda)$.

The corresponding identity contribution is therefore
\begin{equation}
\label{eq:A.identity-localized}
\mathcal{I}(\lambda,\eta)
:=
\vol(M)\int_{\mathbb{R}} \chi_\eta(r-\lambda)\,\mu_{\mathrm{Pl}}(r)\,dr.
\end{equation}

We now extract the leading asymptotic term as $\lambda\to\infty$ in the regime
\begin{equation}
\label{eq:A.window-regime}
\lambda^{-\theta}\le \eta \le 1,
\qquad 0<\theta<\theta_0(\Gamma),
\end{equation}
which is the admissible localization range of the main trace theorem.

\section{Effective volume asymptotics}
\label{sec:A.effective-volume}

The following proposition records the effective volume computation in the localized regime.

\begin{proposition}[Identity term asymptotics]
\label{prop:A.identity-asymptotic}
Let $\chi\in\mathcal{S}(\mathbb{R})$ be even with $\chi(0)=1$.  Then, as $\lambda\to\infty$ uniformly for
$\eta$ in the range \eqref{eq:A.window-regime}, one has
\begin{equation}
\label{eq:A.identity-asymptotic}
\mathcal{I}(\lambda,\eta)
=
\frac{\vol(M)}{2\pi}\,\lambda\,\eta
\int_{\mathbb{R}}\chi(u)\,du
+
O_{\chi}\!\bigl(\vol(M)\,\eta\bigr)
+
O\!\bigl(\vol(M)\,\lambda e^{-2\pi\lambda}\bigr).
\end{equation}
In particular, if $\chi$ is normalized so that $\int_{\mathbb{R}}\chi(u)\,du=2$, then
\begin{equation}
\label{eq:A.identity-asymptotic-simplified}
\mathcal{I}(\lambda,\eta)
=
\frac{\vol(M)}{\pi}\,\lambda\,\eta
+
O_{\chi}\!\bigl(\vol(M)\,\eta\bigr).
\end{equation}
\end{proposition}

\begin{proof}
By the change of variables $r=\lambda+\eta u$ and the definition \eqref{eq:A.plancherel-density},
\eqref{eq:A.identity-localized} becomes
\begin{equation}
\label{eq:A.identity-change}
\mathcal{I}(\lambda,\eta)
=
\vol(M)\,\eta
\int_{\mathbb{R}} \chi(u)\,\mu_{\mathrm{Pl}}(\lambda+\eta u)\,du.
\end{equation}
Using \eqref{eq:A.plancherel-asymptotic} with $r=\lambda+\eta u$, we obtain
\[
\mu_{\mathrm{Pl}}(\lambda+\eta u)
=
\frac{1}{2\pi}(\lambda+\eta u)
+
O\!\bigl((\lambda+\eta|u|)e^{-2\pi(\lambda+\eta|u|)}\bigr).
\]
Substituting into \eqref{eq:A.identity-change} yields
\begin{align*}
\mathcal{I}(\lambda,\eta)
&=
\vol(M)\,\eta
\int_{\mathbb{R}} \chi(u)\left(\frac{1}{2\pi}(\lambda+\eta u)\right)\,du
+
O\!\left(
\vol(M)\,\eta
\int_{\mathbb{R}} |\chi(u)|\,(\lambda+\eta|u|)\,e^{-2\pi(\lambda+\eta|u|)}\,du
\right).
\end{align*}
Since $\chi$ is Schwartz and $\eta\le 1$, the error integral is bounded by
$O(\vol(M)\,\lambda e^{-2\pi\lambda})$.

Expanding the main term gives
\[
\vol(M)\,\eta\cdot \frac{1}{2\pi}
\left(
\lambda \int_{\mathbb{R}}\chi(u)\,du
+
\eta\int_{\mathbb{R}} u\chi(u)\,du
\right).
\]
Since $\chi$ is even, the integral $\int u\chi(u)\,du$ vanishes.  This yields the principal term
\[
\frac{\vol(M)}{2\pi}\,\lambda\,\eta\int_{\mathbb{R}}\chi(u)\,du.
\]
The remaining error term $O_{\chi}(\vol(M)\eta)$ comes from the replacement of
$\mu_{\mathrm{Pl}}(\lambda+\eta u)$ by $\frac{1}{2\pi}(\lambda+\eta u)$ on bounded $u$, together with the rapid decay
of $\chi$ at infinity.  This proves \eqref{eq:A.identity-asymptotic}.
\end{proof}

\subsection{Normalization of the cutoff profile}
\label{subsec:A.cutoff-normalization}

In the main text, it is often convenient to impose the normalization
\begin{equation}
\label{eq:A.cutoff-mass}
\int_{\mathbb{R}}\chi(u)\,du = 2,
\end{equation}
so that the localized window weight approximates the sharp characteristic function of
$[\lambda-\eta,\lambda+\eta]$.

Under \eqref{eq:A.cutoff-mass}, Proposition~\ref{prop:A.identity-asymptotic} yields the universal leading constant
\[
\frac{\vol(M)}{\pi}\,\lambda\eta.
\]
When comparing with the classical local Weyl law expressed in terms of the discrete spectral parameter $r_j$, one
typically absorbs a factor of $2$ depending on whether the window is defined by $|r-\lambda|\le \eta$ or by
$|r-\lambda|\le \eta/2$.  This is a matter of convention and is fixed globally in the main text.

\section{Effective volume and the localized counting function}
\label{sec:A.local-weyl}

Let
\[
N(\lambda,\eta)
:=
\#\{j:\,|r_j-\lambda|\le \eta\},
\]
denote the discrete counting function in a short window.  In the localized trace formula, the identity term
represents the leading asymptotic contribution to the smoothed count
\[
\sum_j \chi_\eta(r_j-\lambda).
\]

The following corollary summarizes the effective volume law that governs the local Weyl constant.

\begin{corollary}[Effective volume constant]
\label{cor:A.effective-volume-constant}
Assume that $\chi$ is even and satisfies \eqref{eq:A.cutoff-mass}.  Then the identity contribution to the localized
trace satisfies
\[
\mathcal{I}(\lambda,\eta)
=
\frac{\vol(M)}{\pi}\,\lambda\,\eta
+
O_{\chi}\!\bigl(\vol(M)\eta\bigr),
\]
uniformly for $\eta$ in the admissible range \eqref{eq:A.window-regime}.
Consequently, the leading constant in the localized Weyl law is given by $\vol(M)/(2\pi)$ per unit spectral length.
\end{corollary}

\begin{proof}
This is immediate from Proposition~\ref{prop:A.identity-asymptotic} under the normalization
$\int \chi = 2$.
\end{proof}

\section{Stability with respect to cutoff modifications}
\label{sec:A.cutoff-stability}

A key feature of the normalized trace formalism developed in Chapter~\ref{chap:08-normalization} is that the
dependence of the trace on the cutoff profile is confined to an explicit identity correction.
For completeness, we record the identity-level stability statement in the present localized setting.

\begin{proposition}[Identity correction under cutoff change]
\label{prop:A.cutoff-correction}
Let $\chi_1,\chi_2\in\mathcal{S}(\mathbb{R})$ be even cutoffs with $\chi_1(0)=\chi_2(0)=1$.  Define the localized
weights $\chi_{j,\eta}(u)=\chi_j(u/\eta)$ and the corresponding identity contributions
$\mathcal{I}_j(\lambda,\eta)$ by \eqref{eq:A.identity-localized}.  Then
\begin{equation}
\label{eq:A.cutoff-correction}
\mathcal{I}_1(\lambda,\eta)-\mathcal{I}_2(\lambda,\eta)
=
\vol(M)\int_{\mathbb{R}}
\bigl(\chi_{1,\eta}(r-\lambda)-\chi_{2,\eta}(r-\lambda)\bigr)\,\mu_{\mathrm{Pl}}(r)\,dr.
\end{equation}
Moreover, as $\lambda\to\infty$ uniformly in \eqref{eq:A.window-regime},
\[
\mathcal{I}_1(\lambda,\eta)-\mathcal{I}_2(\lambda,\eta)
=
\frac{\vol(M)}{2\pi}\,\lambda\,\eta
\int_{\mathbb{R}}(\chi_1(u)-\chi_2(u))\,du
+
O_{\chi_1,\chi_2}\!\bigl(\vol(M)\eta\bigr).
\]
\end{proposition}

\begin{proof}
The identity \eqref{eq:A.cutoff-correction} is immediate from the definition.
The asymptotic follows from the proof of Proposition~\ref{prop:A.identity-asymptotic} applied to $\chi_1-\chi_2$.
\end{proof}

\section{Remarks on normalization conventions}
\label{sec:A.remarks}

The computations above isolate the universal identity constant governing the local Weyl density.  The main term
\[
\frac{\vol(M)}{2\pi}\lambda
\]
should be viewed as the hyperbolic analogue of the Euclidean density in dimension $2$, and it arises entirely from
the Plancherel measure \eqref{eq:A.plancherel-density}.

In the localized trace formula, all other geometric contributions (hyperbolic conjugacy classes, parabolic terms, and
continuous spectrum) are subordinate in the regime $\eta\ge \lambda^{-\theta}$, and their influence is quantified in
the consolidated remainder ledger of Chapter~\ref{chap:08-normalization}.

\section*{Bibliographic notes}

The identity term computation and the Plancherel density for $\mathbb{H}$ are classical.
See \cite{Hejhal1983,Buser1992,Iwaniec2002} for the Selberg trace formula and the derivation of the identity
contribution.  The local Weyl constant and its connection to the Plancherel measure may be found in standard
treatments of spectral asymptotics on locally symmetric spaces.

% ======================================================================
% End of Appendix A
% ======================================================================
